\documentclass{../../note}

\usepackage{amsthm}
\usepackage{pgfplots}
\pgfplotsset{compat=1.18}
\usepackage{xcolor}
\usepackage{tikz}
\usetikzlibrary{shapes,arrows,positioning,fit,calc,matrix,decorations.pathreplacing}
\usepackage{algorithm}
\usepackage{algpseudocode}
\usepackage{listings}
\usepackage{booktabs}
\usepackage{array}
\usepackage{enumitem}

\title{HUFE 数据库原理模拟试卷 B 参考答案}
\author{isomo}
\setlength{\parindent}{0em}

\begin{document}

\fontsize{12pt}{14.4pt}\selectfont
\maketitle

\section*{一、单项选择题}

1.A\quad 2.A\quad 3.C\quad 4.B\quad 5.B\quad 6.B\quad 7.B\quad 8.A\quad 9.C\quad 10.C\quad 11.B\quad 12.B\quad 13.B\quad 14.B\quad 15.A

\section*{二、填空题}

\begin{enumerate}
\item DBMS
\item 除法
\item 数据控制语言
\item 候选码
\item 动态转储
\end{enumerate}

\section*{三、判断题}

1.$\surd$\quad 2.$\surd$\quad 3.$\times$\quad 4.$\surd$\quad 5.$\times$\quad 6.$\times$\quad 7.$\surd$\quad 8.$\surd$\quad 9.$\surd$\quad 10.$\surd$

\section*{四、简答题}

\begin{enumerate}[itemsep=0.8em]
\item 实体完整性要求主码不能为空且唯一；参照完整性要求外码的取值要么为空，要么等于被参照关系中的某个主码值。
\item 事务是数据库中一个不可分割的逻辑工作单位。ACID 分别是原子性、一致性、隔离性、持久性。
\end{enumerate}

\section*{五、综合应用题}

\subsection*{1. SQL 与规范化题}

\begin{enumerate}[label=(\arabic*), itemsep=0.8em]
\item 基本函数依赖：\newline
订单号 $\rightarrow$ 客户号\newline
客户号 $\rightarrow$ 客户名\newline
商品号 $\rightarrow$ 商品名，单价\newline
(订单号，商品号) $\rightarrow$ 数量
\item 候选码：(订单号，商品号)
\item 最高只达到 1NF。因为存在非主属性对候选码的部分依赖，如订单号 $\rightarrow$ 客户号，商品号 $\rightarrow$ 商品名、单价。
\item 可分解为：\newline
订单(\underline{订单号}, 客户号)\newline
客户(\underline{客户号}, 客户名)\newline
商品(\underline{商品号}, 商品名, 单价)\newline
订单明细(\underline{订单号}, \underline{商品号}, 数量)
\end{enumerate}

\subsection*{2. 查询与恢复题}

\begin{enumerate}[label=(\arabic*), itemsep=0.8em]
\item
\begin{lstlisting}[language=SQL]
SELECT Sno, AVG(Grade) AS AvgGrade
FROM SC
GROUP BY Sno
HAVING AVG(Grade) > 80;
\end{lstlisting}

\item
\begin{lstlisting}[language=SQL]
SELECT S.Sname
FROM Student S, SC
WHERE S.Sno = SC.Sno
GROUP BY S.Sno, S.Sname
HAVING COUNT(*) >= 2;
\end{lstlisting}

\item 系统故障恢复通常先从检查点或最近一致状态开始，根据日志文件对已提交事务进行 REDO，对未完成事务进行 UNDO，使数据库恢复到一致状态。
\end{enumerate}

\end{document}
